\SourceLeaf{026}
\Article{14}
\textsc{Lemma.} \textit{If the quantity $r$ and the angle $\varphi$ are determined so that the equations}
\[
\begin{aligned}
&r^m\cos m\varphi+Ar^{m-1}\cos(m-1)\varphi+Br^{m-2}\cos(m-2)\varphi+\text{etc.}\\
&\hspace{3em}+Krr\cos2\varphi+Lr\cos\varphi+M=0\quad.\quad.\quad[1]\\
&r^m\sin m\varphi+Ar^{m-1}\sin(m-1)\varphi+Br^{m-2}\sin(m-2)\varphi+\text{etc.}\\
&\hspace{3em}+Krr\sin2\varphi+Lr\sin\varphi=0\quad.\quad.\quad[2]
\end{aligned}
\]
\textit{hold, the polynomial $x^m+Ax^{m-1}+Bx^{m-2}+\text{etc.}+Kxx+Lx+M=X$ will be divisible by the quadratic factor $xx-2\cos\varphi\mathbin{.}rx+rr$, provided we do not have $r\sin\varphi$ $=0$; but if $r\sin\varphi=0$, the same polynomial will be divisible by the linear factor $x-r\cos\varphi$.}

\textit{Proof.} I. By the preceding article, all the following quantities will be divisible by $xx-2\cos\varphi\mathbin{.}rx+rr$:
{\small
\[
\setlength{\arraycolsep}{2pt}
\begin{array}{rll}
\sin\varphi\mathbin{.}rx^m&-\sin m\varphi\mathbin{.}r^m x&+\sin(m-1)\varphi\mathbin{.}r^{m+1}\\
A\sin\varphi\mathbin{.}rx^{m-1}&-A\sin(m-1)\varphi\mathbin{.}r^{m-1}x&+A\sin(m-2)\varphi\mathbin{.}r^m\\
B\sin\varphi\mathbin{.}rx^{m-2}&-B\sin(m-2)\varphi\mathbin{.}r^{m-2}x&+B\sin(m-3)\varphi\mathbin{.}r^{m-1}\\
&\text{etc.}&\text{etc.}\\
K\sin\varphi\mathbin{.}rxx&-K\sin2\varphi\mathbin{.}rrx&+K\sin\varphi\mathbin{.}r^3\\
L\sin\varphi\mathbin{.}rx&-L\sin\varphi\mathbin{.}rx&\qquad *\\
M\sin\varphi\mathbin{.}r&\qquad *&+M\sin(-\varphi)\mathbin{.}r
\end{array}
\]
}
Their sum will therefore also be divisible by $xx-2\cos\varphi\mathbin{.}rx+rr$. Now the first parts of these quantities form the sum $\sin\varphi\mathbin{.}rX$; the second parts, added together, give $0$ by $[2]$; and it is easy to see that the sum of the third parts also vanishes if we multiply $[1]$ by $\sin\varphi$ and $[2]$ by $\cos\varphi$, and subtract the former product from the latter. It follows that the polynomial $\sin\varphi\mathbin{.}rX$ is divisible by $xx-2\cos\varphi\mathbin{.}rx+rr$, and hence, unless $r\sin\varphi=0$, so is the polynomial $X$. Q. E. P.

II. But if $r\sin\varphi=0$, either $r=0$ or $\sin\varphi=0$. In the first case, $M=0$ by $[1]$, and thus $X$ is divisible by $x$, or by $x-r\cos\varphi$; in the second case, $\cos\varphi=\pm1$, $\cos2\varphi=+1$, $\cos3\varphi=\pm1$, and in general $\cos n\varphi=\cos\varphi^n$. Thus, by $[1]$, we obtain $X=0$ on setting $x=r\cos\varphi$, and consequently the polynomial $X$ will be divisible by $x-r\cos\varphi$. Q. E. S.

\SourceLeaf{027}
\Article{15}
The preceding theorem is usually proved with the aid of imaginary quantities; see \textsc{Euler}, \textit{Introd. in Anal. Inf.}, Vol. I, p. 110. I thought it worth showing how it can be obtained just as easily without their aid. It is now clear that proving our theorem requires nothing more than showing the following: \textit{Given any polynomial $X$ of the form $x^m+Ax^{m-1}+Bx^{m-2}+\text{etc.}+Lx+M$, $r$ and $\varphi$ can be determined so that equations $[1]$ and $[2]$ hold.} For it follows that $X$ has a real factor of the first or second degree; division will necessarily produce a real quotient of lower degree, which, for the same reason, will also have a factor of the first or second degree. By continuing this operation, $X$ will eventually be resolved into real linear or quadratic factors. The purpose of the following investigations is therefore to prove that theorem.

\Article{16}
Imagine a fixed infinite plane (the plane of the plate, Fig. 1), and in it a fixed infinite straight line $GC$ passing through a fixed point $C$. Take some length as the unit, so that all lengths of straight lines can be expressed by numbers. At any point $P$ of the plane whose distance from the centre $C$ is $r$ and whose angle is $GCP=\varphi$, erect a perpendicular equal to the value of the expression
\[
r^m\sin m\varphi+Ar^{m-1}\sin(m-1)\varphi+\text{etc.}+Lr\sin\varphi
\]
which, for brevity, I shall always denote by $T$ in what follows. I always regard the distance $r$ as positive; for points lying on the other side of the axis, the angle $\varphi$ must be regarded either as greater than two right angles or as negative (which here amounts to the same thing). The endpoints of these perpendiculars (which must be taken above the plane for a positive value of $T$, below it for a negative value, and in the plane itself for a vanishing value) will lie on a continuous curved surface extending infinitely in every direction, which, for brevity, I shall call the \textit{first surface} in what follows. In exactly the same way, refer another surface to the same plane and centre and the same axis, with its height above any point of the plane equal to
\[
r^m\cos m\varphi+Ar^{m-1}\cos(m-1)\varphi+\text{etc.}+Lr\cos\varphi+M
\]
an expression which, for brevity, I shall always denote by $U$. I shall distinguish this surface, which will likewise be continuous and extend infinitely in every direction, from the former by calling it the
\SourceLeaf{028}
\textit{second surface}. It is then clear that the whole task is to show that there is at least one point lying in the plane, in the first surface and in the second surface at the same time.

\Article{17}
It is easy to see that the first surface lies partly above and partly below the plane. For the distance $r$ from the centre can clearly be taken so large that the remaining terms in $T$ vanish in comparison with the first term $r^m\sin m\varphi$; and, by choosing the angle $\varphi$ appropriately, this term can be made either positive or negative. The fixed plane will therefore necessarily be cut by the first surface. I shall call this intersection of the plane with the first surface the \textit{first line}; it will thus be determined by the equation $T=0$. For the same reason, the plane will be cut by the second surface; the intersection will form a curve determined by the equation $U=0$, which I shall call the \textit{second line}. Strictly speaking, each curve will consist of several branches, which may be entirely separate, though each will be a continuous line. Indeed, the first line will always be what is called a compound curve, and the axis $GC$ must be regarded as part of this curve: whatever value is assigned to $r$, we always have $T$ $=0$ whenever $\varphi$ $=0$ or $=180^{\circ}$. But it is better to regard the collection of all branches passing through all points where $T=0$ as a single curve (following the usage generally accepted in higher geometry), and likewise all branches passing through all points where $U=0$. It is now clear that the task has been reduced to showing that there is at least one point in the plane where some branch of the first line is cut by a branch of the second line. For this purpose, we must examine the nature of these lines more closely.

\Article{18}
First of all, I observe that both curves are algebraic and, when referred to orthogonal coordinates, of degree $m$. For if the origin of abscissas is taken at $C$, the abscissas $x$ towards $G$, and the ordinates $y$ towards $P$, we have $x=r\cos\varphi$, $y=r\sin\varphi$, and thus in general, whatever $n$ may be,
\[
\begin{aligned}
r^n\sin n\varphi={}&nx^{n-1}y-\frac{n\mathbin{.}n-1\mathbin{.}n-2}{1\mathbin{.}2\mathbin{.}3}x^{n-3}y^3\\
&+\frac{n\mathbin{.}\mathbin{.}\mathbin{.}n-4}{1\mathbin{.}\mathbin{.}\mathbin{.}\mathbin{.}5}x^{n-5}y^5-\text{etc.},\\
r^n\cos n\varphi={}&x^n-\frac{n\mathbin{.}n-1}{1\mathbin{.}2}x^{n-2}yy\\
&+\frac{n\mathbin{.}n-1\mathbin{.}n-2\mathbin{.}n-3}{1\mathbin{.}2\mathbin{.}3\mathbin{.}4}x^{n-4}y^4-\text{etc.}
\end{aligned}
\]
Both $T$ and $U$ will therefore consist of several terms of the form $ax^{\alpha}y^{\beta}$,
\SourceLeaf{029}
where $\alpha$ and $\beta$ denote positive integers whose sum, at its greatest, is $=m$. It can also readily be foreseen that all terms of $T$ contain the factor $y$, and thus that the first line is, strictly speaking, composed of a straight line (whose equation is $y=0$) and a curve of degree $m-1$; but we need not take this distinction into account here.

It will be more important to investigate whether the first and second lines have arms extending to infinity, and how many and of what kind. At an infinite distance from the point $C$, the first line, whose equation is $\sin m\varphi+\frac{A}{r}\sin(m-1)\varphi+\frac{B}{rr}\sin(m-2)\varphi\ \text{etc.}=0$, will coincide with the line whose equation is $\sin m\varphi=0$. This represents $m$ straight lines intersecting at the point $C$, the first of which is the axis $GCG'$, while the others are inclined to it at angles of $\frac{1}{m}180$, $\frac{2}{m}180$, $\frac{3}{m}180$ etc. degrees. The first line therefore has $2m$ infinite branches, which divide the circumference of a circle drawn with infinite radius into $2m$ equal parts: the first branch cuts the circumference at the intersection of the circle and the axis, the second at a distance of $\frac{1}{m}180^{\circ}$, the third at a distance of $\frac{2}{m}180^{\circ}$, etc. In the same way, at an infinite distance from the centre the second line will have an asymptote expressed by the equation $\cos m\varphi=0$. This is a collection of $m$ straight lines likewise intersecting at the point $C$ at equal angles, but in such a way that the first forms an angle of $\frac{1}{m}90^{\circ}$ with the axis $CG$, the second an angle of $\frac{3}{m}90^{\circ}$, the third an angle of $\frac{5}{m}90^{\circ}$, etc. The second line will therefore also have $2m$ infinite branches, each occupying the middle position between two neighbouring branches of the first line, so that they cut the circumference of a circle drawn with an infinitely large radius at points whose angular distances from the axis are $\frac{1}{m}90^{\circ}$, $\frac{3}{m}90^{\circ}$, $\frac{5}{m}90^{\circ}$, etc. It is also clear that the axis itself always forms two infinite branches of the first line, namely the first and the one in position $m+1$. This arrangement of the branches is shown very clearly in Fig. 2, drawn for the case $m=4$; the branches of the second line are dotted to distinguish them from those of the first line, and the same applies to the fourth figure\HistoricalNote{8}. --- Since these conclusions are of the greatest importance, however, and infinitely large quantities may trouble some readers, I shall also show in the next article how to obtain them without the aid of infinities.
